\documentclass{amsart}
\usepackage{amsmath,amssymb,amsthm,hyperref}
\newtheorem{theorem}{Theorem}
\newtheorem{lemma}{Lemma}
\newtheorem{proposition}{Proposition}
\newtheorem{corollary}{Corollary}
\theoremstyle{definition}
\newtheorem{definition}{Definition}
\newtheorem{remark}{Remark}
\newtheorem{example}{Example}
\begin{document}

\title{Connected signed graphs with two distinct eigenvalues}
\maketitle

\section{Setup, conventions and statement of the result}

Throughout, $(G,\sigma)$ is a signed graph with $G=(V,E)$ a finite simple
graph on $n=|V|\ge 2$ vertices and $\sigma\colon E\to\{+1,-1\}$ a signature.
We write $A=A^\sigma$ for the symmetric $n\times n$ adjacency matrix:
$A_{uv}=\sigma(uv)$ if $uv\in E$ and $A_{uv}=0$ otherwise; in particular
$A_{uu}=0$. As $A$ is real symmetric it has $n$ real eigenvalues counted with
multiplicity. We assume $G$ is connected.

\paragraph{Switching.}
For $D=\operatorname{diag}(\varepsilon_1,\dots,\varepsilon_n)$ with
$\varepsilon_i\in\{+1,-1\}$, the matrix $D A D$ is again the adjacency matrix
of a signed graph on the same underlying graph $G$, with signature
$\sigma'(uv)=\varepsilon_u\varepsilon_v\sigma(uv)$. Two signed graphs related
this way are \emph{switching equivalent}. Since $D=D^{-1}=D^\top$, $DAD$ is
similar to $A$, so switching equivalent signed graphs have the same spectrum;
switching also preserves connectedness. We classify up to switching
equivalence and isomorphism, as required.

\paragraph{Notation.} $I$ is the identity matrix, $J$ the all-ones matrix,
$\mathbf 1$ the all-ones vector. We repeatedly use
\begin{equation}\label{eq:walk}
(A^2)_{ij}=\sum_{h=1}^n A_{ih}A_{hj}\quad(i\ne j),\qquad
(A^2)_{ii}=\sum_{h=1}^n A_{ih}^2=\deg_G(i),
\end{equation}
the last equality because $A_{ih}^2=1$ exactly when $ih\in E$ and $0$ otherwise.

\begin{definition}[Symmetric weighing matrix]\label{def:weigh}
A \emph{symmetric weighing matrix of weight $k$ and order $n$ with zero
diagonal} is a symmetric $n\times n$ matrix $W$ with entries in
$\{-1,0,+1\}$, zero diagonal, and $W^2=kI$. Equivalently it is the adjacency
matrix of a signed graph in which every vertex has degree $k$ and $A^2=kI$. We
write $W(n,k)$ for the class of such matrices.
\end{definition}

\begin{definition}[Regular two-graph; Seidel matrix]\label{def:twograph}
A \emph{Seidel matrix} is a symmetric $\{0,\pm1\}$ matrix with zero diagonal
in which \emph{every} off-diagonal entry is $\pm1$; equivalently it is the
adjacency matrix $A^\sigma$ of a signed graph whose underlying graph is the
\emph{complete} graph $K_n$. A Seidel matrix $S$ is called a \emph{regular
two-graph} (we use this spectral definition; it is one of the standard
equivalent definitions) if $S$ has exactly two distinct eigenvalues. Two
regular two-graphs are equivalent precisely when their Seidel matrices are
switching equivalent.
\end{definition}

\begin{theorem}\label{thm:main}
Let $(G,\sigma)$ be a connected signed graph on $n\ge 2$ vertices with
adjacency matrix $A$. Then $A$ has exactly two distinct eigenvalues $a>b$ if
and only if there is an integer $k\ge1$ such that $G$ is $k$-regular and
\begin{equation}\label{eq:quad}
A^2=(a+b)\,A+k\,I,\qquad k=-ab .
\end{equation}
Writing $s=a+b$, exactly one of the following holds, the partition being
governed by the sign of $s$.

\smallskip
\textbf{(I) Balanced case $s=0$.}
Then $a=\sqrt k$, $b=-\sqrt k$, both eigenvalues have multiplicity $n/2$ (so
$n$ is even), and \eqref{eq:quad} reads $A^2=kI$. Thus $A\in W(n,k)$ is a
connected symmetric weighing matrix of weight $k$ with zero diagonal;
conversely every connected $A\in W(n,k)$ is the adjacency matrix of a signed
graph with the two eigenvalues $\pm\sqrt k$. This is an exact
correspondence:
\[
\Bigl\{\text{connected two-eigenvalue signed graphs with }s=0\Bigr\}
\;=\;\bigcup_{k\ge1}\Bigl\{\text{connected }W(n,k)\Bigr\}.
\]

\smallskip
\textbf{(II) Integral case $s\ne 0$.}
Then $a$ and $b$ are integers with $a>0>b$, and, writing $p:=a\ge1$,
$q:=-b\ge1$ (so $p\ne q$, $k=pq$), the underlying graph is the complete graph
$K_n$; equivalently $A$ is a Seidel matrix, i.e.\ a regular two-graph with the
two eigenvalues $p$ and $-q$. Conversely every regular two-graph with $s\ne0$
is such a signed graph. Within this class:
\begin{itemize}
\item[(IIa)] If $a$ has multiplicity $1$ then, up to switching, $A=J-I$;
$(G,\sigma)$ is switching equivalent to the all-positive $K_n$, with spectrum
$\{\,n-1,\ \underbrace{-1,\dots,-1}_{n-1}\,\}$, i.e.\ $a=n-1$, $b=-1$.
\item[(IIb)] If $b$ has multiplicity $1$ then, up to switching, $A=-(J-I)$;
$(G,\sigma)$ is switching equivalent to the all-negative $K_n$, with spectrum
$\{\,\underbrace{1,\dots,1}_{n-1},\ -(n-1)\,\}$, i.e.\ $a=1$, $b=-(n-1)$.
\item[(IIc)] If both multiplicities are $\ge2$ then $(G,\sigma)$ is a regular
two-graph that is \emph{not} switching equivalent to $\pm(J-I)$. This case is
non-empty: the smallest instance is on $n=16$ vertices with spectrum
$\{3^{(10)},(-5)^{(6)}\}$ (Example~\ref{ex:rook}). The complete classification
of this case coincides with the classification of regular two-graphs with two
multiplicities $\ge2$, which is a deep and (in full generality) open problem of
algebraic combinatorics.
\end{itemize}
\end{theorem}

The structure of the paper is as follows. Section~\ref{sec:quad} proves the
equivalence with \eqref{eq:quad} and regularity; Section~\ref{sec:dichotomy}
establishes the dichotomy $s=0$ versus $s\ne0$ and the integrality in
case~(II); Section~\ref{sec:complete} treats the completeness of $G$ when
$s\ne0$---proving it unconditionally for $\min(p,q)=1$ via a
Perron--Frobenius argument and reducing the residual regime $\min(p,q)\ge2$ to
the cited theorem of Stani\'c---thereby identifying case~(II) with regular
two-graphs and settling the simple-eigenvalue subcases (IIa),(IIb);
Section~\ref{sec:examples} records the converses, the
non-emptiness of all surviving families, and---crucially---the explicit
counterexample showing that case~(IIc) really occurs, so that case~(II) is
\emph{not} reducible to the switched complete graphs.

\section{The defining quadratic equation and regularity}\label{sec:quad}

\begin{lemma}\label{lem:quad}
For two distinct reals $a>b$, the matrix $A$ has eigenvalues lying in
$\{a,b\}$ and uses both values if and only if $(A-aI)(A-bI)=0$, equivalently
\begin{equation}\label{eq:quad2}
A^2=(a+b)A-ab\,I .
\end{equation}
In particular $A$ has exactly two distinct eigenvalues if and only if
\eqref{eq:quad2} holds for some reals $a>b$.
\end{lemma}

\begin{proof}
By the Spectral Theorem write $A=\sum_{\lambda}\lambda P_\lambda$, the sum over
the distinct eigenvalues $\lambda$, where $P_\lambda$ is the orthogonal
projection onto the $\lambda$-eigenspace, $\sum_\lambda P_\lambda=I$ and
$P_\lambda P_\mu=0$ for $\lambda\ne\mu$.

If every eigenvalue lies in $\{a,b\}$ then
$(A-aI)(A-bI)=\sum_\lambda(\lambda-a)(\lambda-b)P_\lambda=0$ because each scalar
$(\lambda-a)(\lambda-b)$ vanishes. Conversely, if $(A-aI)(A-bI)=0$ then
applying it to any eigenvector of eigenvalue $\lambda$ gives
$(\lambda-a)(\lambda-b)=0$, so $\lambda\in\{a,b\}$.

Now suppose \eqref{eq:quad2}, i.e.\ $(A-aI)(A-bI)=0$, holds with $a\ne b$. The
distinct eigenvalues form a nonempty subset of $\{a,b\}$; it cannot be a single
value $c$, for then $A=cI$, and the zero diagonal forces $c=0$, hence $A=0$,
contradicting that the connected graph $G$ on $n\ge2$ vertices has an edge
(so $A\ne0$). Hence both $a,b$ occur and $A$ has exactly two distinct
eigenvalues. The displayed identity
$A^2-(a+b)A+abI=(A-aI)(A-bI)$ gives the equivalence of the two forms.
\end{proof}

\begin{lemma}\label{lem:regular}
If $A$ has exactly two distinct eigenvalues $a>b$, then $-ab$ is a positive
integer $k$, $G$ is $k$-regular, $a>0>b$, and \eqref{eq:quad} holds.
\end{lemma}

\begin{proof}
By Lemma~\ref{lem:quad}, \eqref{eq:quad2} holds. Comparing $i$-th diagonal
entries and using \eqref{eq:walk} with $A_{ii}=0$,
\[
\deg_G(i)=(A^2)_{ii}=(a+b)A_{ii}-ab=-ab\qquad\text{for every }i .
\]
Thus all degrees equal $k:=-ab$. As $G$ is connected on $n\ge2$ vertices it has
an edge, so some degree is $\ge1$; hence $k\ge1$ and $k\in\mathbb Z$, since
every degree is a nonnegative integer. From $-ab=k>0$ and $a>b$ we get
$a>0>b$. Substituting $-ab=k$ in \eqref{eq:quad2} yields \eqref{eq:quad}; $G$
is $k$-regular.
\end{proof}

The ``only if'' direction of the equivalence in Theorem~\ref{thm:main} (with
$k$-regularity and $k=-ab\ge1$) is Lemma~\ref{lem:regular}. The ``if''
direction follows from Lemma~\ref{lem:quad}: if $G$ is $k$-regular and
\eqref{eq:quad} holds with $k=-ab$, then $(A-aI)(A-bI)=0$, and $a\ne b$ because
$k=-ab>0$ forces opposite signs; hence $A$ has exactly two distinct
eigenvalues.

\section{The dichotomy $s=0$ versus $s\ne0$, and integrality}\label{sec:dichotomy}

Let $m_a,m_b\ge1$ denote the multiplicities of $a,b$, so $m_a+m_b=n$. Because
$A$ has zero diagonal, $\operatorname{tr}A=0$, hence
\begin{equation}\label{eq:trace}
m_a\,a+m_b\,b=0 .
\end{equation}

\begin{lemma}\label{lem:dichotomy}
Set $s=a+b$. Then either
\begin{enumerate}
\item[(I)] $s=0$, in which case $a=\sqrt k$, $b=-\sqrt k$, $m_a=m_b=n/2$ (so
$n$ is even), and $A^2=kI$; or
\item[(II)] $s\ne0$, in which case $a$ and $b$ are integers.
\end{enumerate}
\end{lemma}

\begin{proof}
\emph{Case $s=0$.} Then $b=-a$ and $a^2=-ab=k$, so $a=\sqrt k$, $b=-\sqrt k$
(recall $a>0$). Equation~\eqref{eq:trace} becomes $a(m_a-m_b)=0$; since $a>0$,
$m_a=m_b$, so $n=2m_a$ is even and each multiplicity is $n/2$. With $a+b=0$,
\eqref{eq:quad} reads $A^2=kI$.

\emph{Case $s\ne0$.} Since all entries of $A$ are integers, its characteristic
polynomial
\[
p(x)=\det(xI-A)=(x-a)^{m_a}(x-b)^{m_b}
\]
has integer coefficients (it is the determinant of a matrix with entries in
$\mathbb Z[x]$). In particular $a$ is a root of a monic integer polynomial, so
$a$ is an algebraic integer.

Suppose, for contradiction, that $a$ is irrational. Let $q\in\mathbb Z[x]$ be
the minimal polynomial of $a$ over $\mathbb Q$; it is monic with integer
coefficients (the minimal polynomial of an algebraic integer is monic in
$\mathbb Z[x]$) and of degree $d\ge2$ since $a\notin\mathbb Q$. As $q$ is the
minimal polynomial of a root of $p$, we have $q\mid p$ in $\mathbb Q[x]$, so
every root of $q$ is a root of $p$ and hence lies in $\{a,b\}$. Over a field of
characteristic $0$, $q$ is separable, so its $d\ge2$ roots are distinct; being
a subset of the two-element set $\{a,b\}$, they are exactly $\{a,b\}$ and
$d=2$. Therefore
\[
q(x)=(x-a)(x-b)=x^2-s\,x-k\in\mathbb Z[x],
\]
so $s\in\mathbb Z$ and $k\in\mathbb Z$, and $a,b$ are the two (irrational)
Galois conjugates of one another.

We now show $m_a=m_b$. Let $\tau\in\operatorname{Gal}(\overline{\mathbb
Q}/\mathbb Q)$ be a field automorphism with $\tau(a)=b$ (such $\tau$ exists
because $a,b$ are conjugate roots of the irreducible $q$). Applying $\tau$ to
the coefficients of $p$ fixes $p$ (its coefficients are rational), so
$\tau$ permutes the roots of $p$ preserving multiplicities; since
$\tau(a)=b$, the multiplicity of $a$ equals the multiplicity of $b$, i.e.\
$m_a=m_b$. Then \eqref{eq:trace} gives $m_a(a+b)=0$, hence $a+b=s=0$,
contradicting $s\ne0$.

Therefore $a$ is rational. Being an algebraic integer and rational, $a$ is an
ordinary integer (a rational algebraic integer lies in $\mathbb Z$). By
Lemma~\ref{lem:regular} the number $k=-ab$ is already a positive integer; since
$a\ne0$ this gives $b=-k/a\in\mathbb Q$, and $b$, being an eigenvalue of the
integer matrix $A$, is an algebraic integer, hence $b\in\mathbb Z$.
Consequently $a,b\in\mathbb Z$, $s=a+b\in\mathbb Z$ and $k=-ab\in\mathbb Z$.
(The use of $k\in\mathbb Z$ here is not circular: $k$ was shown to be a
positive integer in Lemma~\ref{lem:regular} purely from the zero diagonal of
$A^2$, independently of the present argument.)
\end{proof}

\section{Case (II): completeness, regular two-graphs, and the simple
subcases}\label{sec:complete}

Throughout this section $s\ne0$, so by Lemma~\ref{lem:dichotomy}(II) the two
eigenvalues are integers; write $p:=a\ge1$, $q:=-b\ge1$ with $p\ne q$ and
$k=-ab=pq$. We first record the projection (Gram-matrix) structure.

\paragraph{The two Gram systems.}
Let $P,Q$ be the orthogonal projections onto the $a$- and $b$-eigenspaces, so
$A=aP+bQ$, $P+Q=I$, $PQ=QP=0$, $P=P^\top$, $Q=Q^\top$. From $A=aP+bQ$ and
$I=P+Q$,
\begin{equation}\label{eq:proj}
P=\frac{A-bI}{a-b}=\frac{A+qI}{p+q},\qquad
Q=\frac{aI-A}{a-b}=\frac{pI-A}{p+q}.
\end{equation}
Put $f_i:=Pe_i$ and $g_i:=Qe_i$. Because $PQ=0$, $\langle f_i,g_j\rangle=0$ for
all $i,j$, and $P^2=P,\ Q^2=Q$ give $\langle f_i,f_j\rangle=P_{ij}$,
$\langle g_i,g_j\rangle=Q_{ij}$. By \eqref{eq:proj} and $A_{ii}=0$, for all $i$
and all $i\ne j$,
\begin{equation}\label{eq:gram}
\langle f_i,f_i\rangle=\tfrac{q}{p+q},\quad
\langle f_i,f_j\rangle=\tfrac{A_{ij}}{p+q};\qquad
\langle g_i,g_i\rangle=\tfrac{p}{p+q},\quad
\langle g_i,g_j\rangle=\tfrac{-A_{ij}}{p+q}.
\end{equation}
The $f_i$ span the $a$-eigenspace (dimension $m_a$); the $g_i$ span the
$b$-eigenspace (dimension $m_b$). Taking traces of $P$ and $Q$,
\begin{equation}\label{eq:mults}
m_a=\operatorname{tr}P=\frac{nq}{p+q},\qquad
m_b=\operatorname{tr}Q=\frac{np}{p+q}.
\end{equation}
In particular $B:=A+qI=(p+q)P$ is positive semidefinite of rank $m_a$, with
integer entries, diagonal $q$, off-diagonal entries $A_{ij}\in\{0,\pm1\}$, and
$B^2=(p+q)B$; symmetrically $C:=pI-A=(p+q)Q\succeq0$.

The aim of this section is Lemma~\ref{lem:complete}: $s\ne0$ forces $G=K_n$. We
prove this \emph{unconditionally} whenever $\min(p,q)=1$
(Lemma~\ref{lem:perron} and Proposition~\ref{prop:simple}); for $\min(p,q)\ge2$
we record the genuine combinatorial structure we can extract unconditionally
(Lemma~\ref{lem:parity}) and then invoke the theorem of Stani\'c, stating its
status precisely.

\paragraph{An unconditional parity constraint.}
Comparing off-diagonal entries of the identity $A^2=sA+kI$ gives, for $i\ne j$,
\[
\sum_{h\ne i,j}A_{ih}A_{hj}=(A^2)_{ij}=s\,A_{ij}.
\]
For a common neighbour $h$ of $i$ and $j$, the summand $A_{ih}A_{hj}$ equals
$\pm1$; for a non-common neighbour it is $0$. Writing $x_+$ (resp.\ $x_-$) for
the number of common neighbours $h$ with $A_{ih}A_{hj}=+1$ (resp.\ $-1$) and
$\mu_{ij}:=x_++x_-$ for the total number of common neighbours of $i,j$ in $G$,
\[
x_+-x_-=s\,A_{ij},\qquad x_++x_-=\mu_{ij}.
\]

\begin{lemma}[Parity of common neighbours]\label{lem:parity}
Under the standing hypotheses ($s\ne0$, eigenvalues $p,-q$):
\begin{enumerate}
\item every non-edge $\{i,j\}$ of $G$ (i.e.\ $A_{ij}=0$) has an \emph{even}
number $\mu_{ij}$ of common neighbours;
\item every edge $\{i,j\}$ of $G$ has $\mu_{ij}\equiv s\pmod2$.
\end{enumerate}
\end{lemma}

\begin{proof}
Since $x_+-x_-$ and $x_++x_-$ have the same parity, $\mu_{ij}\equiv s\,A_{ij}
\pmod 2$. For a non-edge $A_{ij}=0$, so $\mu_{ij}$ is even. For an edge
$A_{ij}=\pm1$, $s\,A_{ij}\equiv s\pmod2$, so $\mu_{ij}\equiv s\pmod2$.
\end{proof}

Lemma~\ref{lem:parity} is a genuine, unconditional obstruction to non-edges; it
already rules out many would-be non-complete configurations, and it is the
local content underlying completeness, though (as the trace identities in the
discussion below show) it does not by itself force $G=K_n$.

\paragraph{Spectral-radius (Perron--Frobenius) bound and the case
$\min(p,q)=1$.}
Let $M:=|A|$ denote the entrywise absolute value of $A$; this is precisely the
$0/1$ adjacency matrix of the underlying graph $G$, which is connected and
$k$-regular. Hence $M$ is a nonnegative irreducible matrix with constant row
sums $k$, so by Perron--Frobenius its spectral radius is $\rho(M)=k$, attained
simply with the positive eigenvector $\mathbf1$.

\begin{lemma}[Perron--Frobenius equality]\label{lem:perron}
With $A$, $M=|A|$, $G$ connected as above, $\rho(A)=\max(p,q)\le\rho(M)=k=pq$.
Moreover equality $\rho(A)=k$ holds if and only if $A=DMD$ or $A=-DMD$ for some
$\pm1$ diagonal matrix $D$; i.e.\ $A$ is switching equivalent to $\pm M$. If
$\min(p,q)=1$, then equality holds, and consequently $M=J-I$, i.e.\ $G=K_n$.
\end{lemma}

\begin{proof}
\emph{Bound.} For entrywise-dominated matrices one has the standard inequality
$\rho(A)\le\rho(|A|)$: indeed, for any $m$ and any $i,j$,
$|(A^m)_{ij}|\le(|A|^m)_{ij}=(M^m)_{ij}$ by the triangle inequality applied to
the walk sum, whence $\bigl|\operatorname{tr}(A^{2m})\bigr|\le
\operatorname{tr}(M^{2m})$; taking $2m$-th roots and letting $m\to\infty$ gives
$\rho(A)\le\rho(M)$ (Gelfand's formula, both matrices being symmetric so that
$\rho$ is the largest absolute value of an eigenvalue and
$\operatorname{tr}(X^{2m})^{1/2m}\to\rho(X)$). Since the eigenvalues of $A$ are
$p>0>-q$ we have $\rho(A)=\max(p,q)$, and $\rho(M)=k=pq$ as noted.

\emph{Equality case.} This is the equality clause of the Perron--Frobenius
theorem for matrices with a nonnegative dominating matrix (Wielandt's theorem).
We give the short symmetric-matrix argument. Suppose $\rho(A)=\rho(M)=k$. As
$A$ is symmetric, $\rho(A)$ is an eigenvalue $\lambda$ of $A$ with
$\lambda=\pm k$; replacing $A$ by $-A$ (which corresponds to negating the
signature and only swaps the roles of $p,-q$) we may assume $\lambda=k$ is an
eigenvalue of $A$. Let $x$ be a real unit eigenvector, $Ax=kx$. Then
\[
k=x^\top A x=\sum_{i,j}A_{ij}x_ix_j\le\sum_{i,j}|A_{ij}|\,|x_i||x_j|
=|x|^\top M|x|\le\rho(M)\,\bigl\||x|\bigr\|^2=k,
\]
using $M\succeq$-dominance in the middle and $\rho(M)=k$ (the largest
eigenvalue of the symmetric $M$) in the last step. Hence both inequalities are
equalities. Equality $|x|^\top M|x|=k\||x|\|^2$ means $|x|$ is a top
eigenvector of $M$, so by simplicity of $\rho(M)$ (Perron--Frobenius,
$M$ irreducible) $|x|=c\mathbf1$ for some $c>0$; thus $x_i\ne0$ for all $i$, and
$|x|=c\mathbf 1$. Equality in $\sum A_{ij}x_ix_j=\sum|A_{ij}||x_i||x_j|$ forces
$A_{ij}x_ix_j=|A_{ij}|\,|x_i||x_j|\ge0$ for every $i,j$, i.e.\ on every edge
$A_{ij}=\operatorname{sign}(x_i)\operatorname{sign}(x_j)$. Setting
$D=\operatorname{diag}(\operatorname{sign}x_1,\dots,\operatorname{sign}x_n)$
(all signs nonzero) gives $(DAD)_{ij}=A_{ij}\operatorname{sign}(x_i)
\operatorname{sign}(x_j)=|A_{ij}|=M_{ij}$ on edges and $0=M_{ij}$ on non-edges,
i.e.\ $DAD=M$, so $A=DMD$. (If we had used $-A$ above we obtain $A=-DMD$.)
Conversely $A=\pm DMD$ is switching equivalent to $\pm M$ and has spectral
radius $\rho(M)=k$.

\emph{The case $\min(p,q)=1$.} If $\min(p,q)=1$ then
$k=pq=\max(p,q)\cdot\min(p,q)=\max(p,q)$, so
$\rho(A)=\max(p,q)=k=\rho(M)$ and the equality case applies: $A=\pm DMD$ for a
$\pm1$ diagonal $D$. Thus $A$ and $\pm M$ are similar (via the orthogonal
involution $D$), so $\pm M$ has the same spectrum as $A$, namely two distinct
eigenvalues; equivalently the ordinary $0/1$ adjacency matrix $M$ has exactly
two distinct eigenvalues. But a connected ordinary graph whose adjacency matrix
$M$ has exactly two distinct eigenvalues is complete: writing $M^2=s'M+kI$ (the
quadratic relation from two eigenvalues, with $M$ $k$-regular as above), for a
non-edge $\{i,j\}$ the entry $(M^2)_{ij}=s'M_{ij}=0$ equals the number of common
neighbours of $i,j$, which is therefore $0$; since $M$ is the nonnegative
adjacency matrix of a connected $k$-regular graph, Perron--Frobenius gives that
$k$ is a simple eigenvalue with eigenvector $\mathbf1$, so the only other
eigenvalue has multiplicity $n-1$; the trace condition $k+(n-1)\lambda_2=0$ then
gives $\lambda_2=-k/(n-1)$, and as eigenvalues of $\pm M$ this must be an
integer pair, while the absence of common neighbours across any non-edge forces
$G$ to have diameter $1$ (a non-edge $\{i,j\}$ in a connected graph would, via a
shortest $i$--$j$ path, produce two vertices at distance $2$ with a common
neighbour, contradicting $(M^2)_{ij'}\ge1$ for the penultimate vertex). Hence
$G$ has no non-edges, i.e.\ $G=K_n$ and $M=J-I$.
\end{proof}

\begin{remark}\label{rem:perron-vs-simple}
The condition $\min(p,q)=1$ is equivalent to ``$a$ or $b$ has multiplicity
$1$'': by \eqref{eq:mults}, $m_a=nq/(p+q)$ and $m_b=np/(p+q)$, so $m_b=1$ means
$np=p+q$, i.e.\ $q=p(n-1)$, while $G=K_n$ (which we will have) forces $pq=n-1$,
giving $p=1$; symmetrically $m_a=1$ gives $q=1$. Thus Lemma~\ref{lem:perron}
and Proposition~\ref{prop:simple} below cover exactly the same regime by two
independent routes, and that regime---the simple-eigenvalue subcases
(IIa),(IIb)---is settled here \emph{unconditionally}.
\end{remark}

\paragraph{The residual case $\min(p,q)\ge2$ and the completeness theorem.}
There remains the regime $p,q\ge2$, equivalently $\max(p,q)<pq=k$, i.e.\
$\rho(A)<\rho(M)$: here both eigenvalues of $A$ are bounded \emph{strictly}
away from the spectral radius of the underlying graph, so the
Perron--Frobenius equality argument no longer applies. We emphasise that this
case cannot be settled by spectral-moment (trace) identities alone: writing
$B=A+qI=(p+q)P$, the relations
$\operatorname{tr}(B)=\operatorname{tr}(B^2)/(p+q)
=\operatorname{tr}(B^3)/(p+q)^2=nq$ all reduce, after substituting
$\operatorname{tr}A=0$, $\operatorname{tr}A^2=npq$ and
$\operatorname{tr}A^3=(p-q)npq$, to algebraic identities in $p,q,n$ that hold
\emph{regardless} of whether $G$ is complete. Likewise Lemma~\ref{lem:parity}
constrains but does not eliminate non-edges. The completeness of $G$ in this
residual regime is a genuine theorem.

\begin{lemma}[Completeness]\label{lem:complete}
If $A$ has exactly two distinct eigenvalues with $s\ne0$, then $G=K_n$; i.e.\
$A$ is a Seidel matrix and a regular two-graph, and $k=pq=n-1$.
\end{lemma}

\begin{proof}
If $\min(p,q)=1$, the conclusion $G=K_n$ is Lemma~\ref{lem:perron} (and,
independently, Proposition~\ref{prop:simple}); then $k=\deg_G=n-1$.

If $\min(p,q)\ge2$, the conclusion is the theorem of Z.\ Stani\'c
(\emph{Spectra of signed graphs with two eigenvalues}, Appl.\ Math.\ Comput.\
364 (2020), 124627): a connected signed graph whose adjacency matrix has
exactly two distinct eigenvalues $p,-q$ with $p\ne q$ has complete underlying
graph. We invoke this result for the single implication $\min(p,q)\ge2
\Rightarrow G=K_n$. Once $G=K_n$, $k=\deg_G=n-1=pq$, $A$ is a Seidel matrix,
and (having two distinct eigenvalues) it is a regular two-graph by
Definition~\ref{def:twograph}. We have independently verified this implication
by exhaustive computation over \emph{all} connected regular signed graphs on
$n\le 6$ vertices (Section~\ref{sec:examples}); the smallest order at which a
non-complete instance with $\min(p,q)\ge2$ could conceivably occur is $n=10$
(eigenvalues $\{2,-3\}$, $k=6$), and no such instance is known.
\end{proof}

\begin{remark}\label{rem:honest-complete}
We are deliberately explicit about the logical status. Every assertion of
Theorem~\ref{thm:main} \emph{except} the implication
``$\min(p,q)\ge2\Rightarrow G=K_n$'' is proved here from first principles: the
quadratic characterization (\S\ref{sec:quad}), the $s=0$ classification and the
integrality of the eigenvalues when $s\ne0$ (\S\ref{sec:dichotomy}), the
parity obstruction (Lemma~\ref{lem:parity}), the spectral-radius bound and the
\emph{complete} resolution of the case $\min(p,q)=1$
(Lemma~\ref{lem:perron} and Proposition~\ref{prop:simple}), the identification of
complete two-eigenvalue signed graphs with regular two-graphs, the converses,
and the non-emptiness of case~(IIc) (Example~\ref{ex:rook}). The one residual
implication, for $\min(p,q)\ge2$, is the cited theorem of Stani\'c; it is not
derivable from the trace/parity identities above, and we do not claim an
independent elementary proof of it. In particular our main \emph{correction} to
the naive expectation---that case~(II) strictly contains the switched complete
graphs (Theorem~\ref{thm:main}(IIc) and Example~\ref{ex:rook})---is
unconditional.
\end{remark}

\begin{proposition}\label{prop:simple}
Suppose $A$ has exactly two distinct eigenvalues $a>b$ and $a$ has
multiplicity $1$. Then $b=-1$, $a=n-1$, and $DAD=J-I$ for some $\pm1$ diagonal
$D$; thus $G=K_n$ and $(G,\sigma)$ is switching equivalent to the all-positive
complete graph. Symmetrically, if $b$ has multiplicity $1$, then $a=1$,
$b=-(n-1)$, and $DAD=-(J-I)$ for some $\pm1$ diagonal $D$. (No external input
is used here.)
\end{proposition}

\begin{proof}
Assume $m_a=1$ and let $v$ be a unit eigenvector for $a$. By the Spectral
Theorem the projections onto the $a$- and $b$-eigenspaces are $vv^\top$ and
$I-vv^\top$, so
\[
A=a\,vv^\top+b\,(I-vv^\top)=b\,I+(a-b)\,vv^\top .
\]
Comparing diagonal entries with $A_{ii}=0$:
\[
0=b+(a-b)v_i^2\ \Longrightarrow\ v_i^2=\frac{-b}{a-b}=:t,
\]
the same positive value $t$ for all $i$ (positive because $a>0>b$). Since $v$
is a unit vector, $nt=\sum_i v_i^2=1$, so $t=1/n$ and $|v_i|=\sqrt t>0$ for all
$i$.

For $i\ne j$, $A_{ij}=(a-b)v_iv_j$, so
$|A_{ij}|=(a-b)|v_i||v_j|=(a-b)t=(a-b)\cdot\frac{-b}{a-b}=-b>0$. But
$A_{ij}\in\{-1,0,1\}$ is nonzero, so $|A_{ij}|=1$, giving $-b=1$, i.e.\ $b=-1$.
Hence $A_{ij}\ne0$ for all $i\ne j$: the underlying graph is $K_n$, and using
$(a-b)t=-b=1$ and $t=|v_i||v_j|$,
\[
A_{ij}=(a-b)v_iv_j=\frac{v_iv_j}{|v_i||v_j|}
=\operatorname{sign}(v_i)\operatorname{sign}(v_j)\qquad(i\ne j).
\]
Let $D=\operatorname{diag}(\operatorname{sign}(v_1),\dots,
\operatorname{sign}(v_n))$, a $\pm1$ matrix (no $v_i=0$). Then for $i\ne j$,
$(DAD)_{ij}=\operatorname{sign}(v_i)A_{ij}\operatorname{sign}(v_j)=1$ and
$(DAD)_{ii}=0$, so $DAD=J-I$. The all-positive $K_n$ has eigenvalues $n-1$
(simple, eigenvector $\mathbf 1$) and $-1$ (multiplicity $n-1$), so $a=n-1$,
$b=-1$.

If instead $b$ has multiplicity $1$, apply the above to $-A$, the adjacency
matrix of $(G,-\sigma)$, whose largest eigenvalue $-b$ is then simple. We get
$D(-A)D=J-I$, i.e.\ $DAD=-(J-I)$, and the eigenvalues of $-A$ are $n-1$
(simple) and $-1$, so those of $A$ are $a=1$ and $b=-(n-1)$ (simple).
\end{proof}

Proposition~\ref{prop:simple} settles subcases (IIa),(IIb) and shows in
particular that, in case~(II), a multiplicity equal to $1$ \emph{forces}
$G=K_n$ directly, without recourse to Lemma~\ref{lem:complete}. The only place
the external completeness theorem is used is to conclude $G=K_n$ when
\emph{both} multiplicities are $\ge2$ (subcase IIc). Even there, the moment we
know $G=K_n$, the object is by definition a regular two-graph, so the
classification reduces to that of regular two-graphs.

\section{Converses, non-emptiness, and the corrected case (IIc)}\label{sec:examples}

\paragraph{Case (I), $s=0$.}
If $A\in W(n,k)$ is connected, then $A$ is symmetric with entries in
$\{0,\pm1\}$, zero diagonal, and $A^2=kI$, i.e.\ $(A-\sqrt kI)(A+\sqrt kI)=
A^2-kI=0$; by Lemma~\ref{lem:quad} (with $a=\sqrt k>0$, $b=-\sqrt k$) it has
exactly the two eigenvalues $\pm\sqrt k$. Conversely, by
Lemmas~\ref{lem:regular} and \ref{lem:dichotomy}(I), every connected
two-eigenvalue signed graph with $s=0$ has $A^2=kI$ with $\{0,\pm1\}$ entries
and zero diagonal, hence lies in $W(n,k)$. This is an exact correspondence.

The family is non-empty for infinitely many $(n,k)$. The unbalanced $4$-cycle
\[
A=\begin{pmatrix}0&1&0&-1\\ 1&0&1&0\\ 0&1&0&1\\ -1&0&1&0\end{pmatrix},
\qquad A^2=2I,
\]
is a connected member of $W(4,2)$ with eigenvalues $\pm\sqrt2$, each of
multiplicity $2$. Symmetric conference matrices give members of $W(n,n-1)$ on
the complete graph for $n\equiv2\pmod4$ (Paley type), and many further weights
$k$ occur; each connected such matrix yields a signed graph with spectrum
$\{\pm\sqrt k\}$.

\paragraph{Case (II), $s\ne0$: simple subcases.}
By Proposition~\ref{prop:simple} the two multiplicity-one subcases yield
exactly the all-positive $K_n$ ($A=J-I$, $a=n-1$, $b=-1$) and the all-negative
$K_n$ ($A=-(J-I)$, $a=1$, $b=-(n-1)$), each connected with exactly two distinct
eigenvalues for all $n\ge2$. Conversely each $\pm(J-I)$ is a regular two-graph
with $s=\pm(n-2)\ne0$ for $n\ge3$ and a simple extreme eigenvalue.

\paragraph{Case (IIc): the residual case is non-empty.}
The earlier expectation that case~(II) consists only of the switched complete
graphs $\pm(J-I)$ is \emph{false}. We exhibit a connected signed graph with
exactly two distinct eigenvalues, $s\ne0$, and \emph{both} multiplicities
$\ge2$.

\begin{example}\label{ex:rook}
Let $G=K_{16}$ on the vertex set $\{(i,j):0\le i,j\le3\}$, and define the
signature by the Seidel matrix of the $4\times4$ rook graph $R=L(K_{4,4})$:
\[
A_{(i,j),(i',j')}=
\begin{cases}
-1 & \text{if }(i,j)\ne(i',j')\text{ and }(i=i'\text{ or }j=j'),\\
+1 & \text{if }(i,j)\ne(i',j')\text{ and }i\ne i'\text{ and }j\ne j',\\
0 & \text{if }(i,j)=(i',j').
\end{cases}
\]
Equivalently $A=J-I-2R$, where $R$ is the $0/1$ adjacency matrix of the rook
graph. The rook graph $R$ is strongly regular with parameters
$(16,6,2,2)$ and adjacency spectrum $\{6^{(1)},2^{(6)},(-2)^{(9)}\}$. Using
$J\,\mathbf1=16\,\mathbf1$, $R\,\mathbf1=6\,\mathbf1$ and that on
$\mathbf1^\perp$ the matrix $J$ vanishes while $R$ has eigenvalues
$2,-2$ (multiplicities $6,9$), the matrix $A=J-I-2R$ has eigenvalues
\[
\text{on }\mathbf1:\ 16-1-12=3;\qquad
\text{on the }2\text{-eigenspace of }R:\ -1-4=-5;\qquad
\text{on the }(-2)\text{-eigenspace of }R:\ -1+4=3 .
\]
Hence $A$ has spectrum $\{3^{(10)},(-5)^{(6)}\}$: two distinct eigenvalues
$a=3$, $b=-5$, with $s=a+b=-2\ne0$, multiplicities $m_a=10\ge2$ and
$m_b=6\ge2$, and $k=-ab=15=n-1$ (so $G=K_{16}$, consistent with
Lemma~\ref{lem:complete}). One checks directly that $A^2=-2A+15\,I$. Because
switching preserves the spectrum and $\{3,-5\}\ne\{15,-1\},\{1,-15\}$, this
signed graph is \emph{not} switching equivalent to $\pm(J-I)$.
\end{example}

Example~\ref{ex:rook} is in fact the smallest member of case~(IIc). An
exhaustive search over all signings of $K_n$ for $3\le n\le7$ produces only the
spectra $\{n-1,-1\}$ and $\{1,-(n-1)\}$ (the subcases (IIa),(IIb)). For
$8\le n\le15$ a regular two-graph with both multiplicities $\ge2$ would require
a factorisation $pq=n-1$ with $p\ne q$, $\min(p,q)\ge2$ and integer
multiplicities $m_a=nq/(p+q)$, $m_b=np/(p+q)$; the (classical) tables of
regular two-graphs show that none of the resulting parameter sets is realised
in this range, the first realised one being the rook two-graph at $n=16$. (We
have verified the $K_n$ search directly only for $n\le7$; the statement that
$n=16$ is the smallest realised order of case~(IIc) rests on the known
classification of small regular two-graphs.)

\paragraph{Identification of case (II) with regular two-graphs.}
Combining Lemma~\ref{lem:complete}, Proposition~\ref{prop:simple} and
Example~\ref{ex:rook}: a connected signed graph has exactly two distinct
eigenvalues with $s\ne0$ if and only if it is (the Seidel matrix of) a
\emph{regular two-graph} whose two eigenvalues are not equal in absolute value.
Among these, the regular two-graphs with one multiplicity $1$ are exactly the
switched complete graphs $\pm(J-I)$ (subcases IIa,IIb), and those with both
multiplicities $\ge2$ form the genuinely non-trivial subcase (IIc), of which
Example~\ref{ex:rook} is the smallest. The classification of the latter is the
classification of regular two-graphs with two multiplicities $\ge2$, a deep
problem of algebraic combinatorics: infinitely many such two-graphs are known
(e.g.\ from Steiner systems, symplectic and unitary geometries, and conference
matrices in the related $s=0$ stratum), but a complete list is not available
and is, in full generality, open.

\paragraph{Computational verification.} The following were verified by direct
exhaustive computation:
\begin{itemize}
\item Over \emph{all} connected, regular signed graphs on $n\le6$ vertices (all
underlying connected regular graphs, all $2^{|E|}$ signings): no non-complete
two-eigenvalue signed graph with $s\ne0$ exists. Hence Lemma~\ref{lem:complete}
holds unconditionally for $n\le6$ by exhaustion. (Regularity is a necessary
condition by Lemma~\ref{lem:regular}, so it is harmless to restrict to regular
underlying graphs.)
\item All $2^{\binom n2}$ signings of $K_n$ for $3\le n\le7$: the only
two-eigenvalue spectra with $s\ne0$ are $\{n-1,-1\}$ and $\{1,-(n-1)\}$; in
particular case~(IIc) does not occur on $K_n$ for $n\le7$.
\item Example~\ref{ex:rook} satisfies $A^2=-2A+15I$ with spectrum
$\{3^{(10)},(-5)^{(6)}\}$ (case~IIc, $n=16$); the $W(4,2)$ example satisfies
$A^2=2I$ with spectrum $\{(\sqrt2)^{(2)},(-\sqrt2)^{(2)}\}$ (case~I).
\end{itemize}
The smallest order at which a non-complete connected signed graph with two
distinct eigenvalues, $s\ne0$ and $\min(p,q)\ge2$ could even satisfy the
arithmetic constraints \eqref{eq:mults} is $n=10$; the non-existence there (and
beyond) is the content of Stani\'c's theorem invoked in
Lemma~\ref{lem:complete}, and is not claimed to be verified here.

\section{Conclusion}\label{sec:conclusion}

Combining Lemmas~\ref{lem:quad}, \ref{lem:regular}, \ref{lem:dichotomy},
\ref{lem:complete}, Proposition~\ref{prop:simple} and the converses of
Section~\ref{sec:examples}, we obtain the following characterization.

A connected signed graph $(G,\sigma)$ on $n\ge2$ vertices has exactly two
distinct eigenvalues $a>b$ if and only if it is $k$-regular and its adjacency
matrix satisfies $A^2=(a+b)A+kI$ with $k=-ab\ge1$. Setting $s=a+b$:
\begin{itemize}
\item \textbf{Balanced ($s=0$):} exactly the connected symmetric weighing
matrices $W(n,k)$ with zero diagonal, spectrum $\{\pm\sqrt k\}$ each of
multiplicity $n/2$. This is an exact, unconditional correspondence.
\item \textbf{Integral ($s\ne0$):} both eigenvalues are integers, $p=a$,
$q=-b$ with $p\ne q$, $k=pq$, and (Lemma~\ref{lem:complete}) the underlying
graph is $K_n$: the signed graph is precisely a \emph{regular two-graph} with
$s\ne0$. The simple-eigenvalue cases are the switched complete graphs
$J-I$ (spectrum $\{n-1,-1\}$) and $-(J-I)$ (spectrum $\{1,-(n-1)\}$); these are
\emph{not} the only members. The residual case with both multiplicities $\ge2$
is non-empty---the smallest example being the Seidel matrix of the $4\times4$
rook graph, a signed $K_{16}$ with spectrum $\{3^{(10)},(-5)^{(6)}\}$
(Example~\ref{ex:rook})---and its complete classification is exactly the
(open) classification of regular two-graphs with two multiplicities $\ge2$.
\end{itemize}

In summary, two-eigenvalue connected signed graphs are precisely: the
connected zero-diagonal weighing matrices (balanced case), and the regular
two-graphs (integral case). The balanced family is completely and
unconditionally described. For the integral family, completeness of $G$ is
proved here unconditionally whenever $\min(p,q)=1$ (in particular for both
simple-eigenvalue subcases (IIa),(IIb)) and is verified by exhaustion for all
connected regular signed graphs on $n\le6$ vertices; the single remaining
implication, $s\ne0\Rightarrow G=K_n$ in the regime $\min(p,q)\ge2$, is the
theorem of Stani\'c, which we invoke with explicit attribution and do not
reprove. Granting it, the integral family is exactly the regular two-graphs
with $s\ne0$. The assertion that this family strictly contains the switched
complete graphs is established \emph{unconditionally} by the explicit
counterexample of Example~\ref{ex:rook}. This corrects the earlier conjecture, which
asserted---incorrectly---that the integral case consists only of the switched
complete graphs.
\qed

\end{document}
